\begin{table}[!htb]\centering
\caption{The first years of the alternative transitions}\label{tab:alttrans}
\begingroup\small
\sbox{\ntbox}{\begin{tabular}{@{}lcccccc@{}}
\toprule
 & & \multicolumn{2}{c}{Homelessness} & \multicolumn{2}{c}{Children in a shelter} & Newborn's \\
\cmidrule(lr){3-4}\cmidrule(lr){5-6}
 & $p_{\mathrm{det}}$ & Largest rise & Periods above & Largest rise & Periods above & gain \\
\midrule
Main transition & 20 & none & none & none & none & 1.45 \\
\addlinespace
\multicolumn{7}{@{}l}{No exemption} \\
 & 100 & 13.6 (2) & 1--15 & 10.8 (2) & 2--11 & 6.57 \\
 & 20 & 5.4 (7) & 2--17 & 3.1 (5) & 2--11 & 3.09 \\
\addlinespace
\multicolumn{7}{@{}l}{Announced five periods ahead} \\
 & 100 & 7.3 (9) & 7--16 & 3.6 (9) & 7--12 & 4.54 \\
 & 20 & 3.4 (12) & 7--19 & 0.9 (11) & 8--13 & 1.92 \\
\addlinespace
\multicolumn{7}{@{}l}{Sobriety while in treatment exempt for the first five periods} \\
 & 100 & 5.2 (7) & 7--14 & 1.7 (7) & 7--10 & 6.24 \\
 & 20 & 1.6 (11) & 9--16 & none & none & 2.87 \\
\bottomrule
\end{tabular}}
\usebox{\ntbox}
\tabnote{Notes: Each row is a transition from the Housing First benchmark to Treatment First with the drug price and $\pi_V$ clearing every period, against Housing First's steady state. Main transition: Section~\ref{sec:ftrans}, where a household with a SUD keeps its voucher by entering treatment and not using. In the other rows it cannot stop using, and the disorder fails the screen even in treatment: from the reform (no exemption), from period six with Housing First until then (announced five periods ahead), or from period six with treatment alone passing it until then (sobriety while in treatment exempt for the first five periods). $p_{\mathrm{det}}$ is the detection rate, in percent. Largest rise: the largest percent increase over the first fifteen periods, with its period in brackets. Periods above: the periods in which it is above Housing First. Children in a shelter is the homelessness of households aged 26 to 50. Newborn's gain: the CEV of children born at the reform, percent of lifetime consumption. When every user is caught, no drug price clears the long-run drug market, because drug use jumps as the price moves through the level that would clear it; those transitions end at the long-run solution that comes closest.}
\endgroup
\end{table}
